\section{Flags, incidence, and selected-tuple entropy}
\label{base:section}

Fix an integer $d\ge5$ and $0<\eta<1/10$. Throughout the construction,
$q$ is a sufficiently large prime, and
\begin{equation}
 \sigma=\log q,\qquad N=\lfloor q^d\sigma\rfloor,
 \qquad k=\lfloor q\sigma^{1+\eta}\rfloor.
 \label{base:stream-parameters}
\end{equation}
Constants may depend on $d,\eta$ and on fixed retention and budget
constants, but not on $q$. In particular, the constants implicit in
$\ell=\Theta(k)$ do not change with $q$.
We first construct a graph that always excludes $K_{d+1}$. The remaining
task will be to exclude a consistent $k$-tuple of its flags. The entropy
bound below identifies the amount of information that any selected tuple
must retain, even when its selection is highly biased.

\subsection{Projective incidence}

The following point--hyperplane calculation is the bipartite form of the
projective orthogonality estimate used by Alon and Krivelevich
\cite{AlonKrivelevich1997}; we include the weighted identity needed later.

For $j\ge1$, let $V_j=\F_q^{j+1}$ and let $V_j^*$ be its dual vector
space. The two projective spaces $\PG(V_j)$ and $\PG(V_j^*)$ consist of
one-dimensional vector subspaces; both have size
\[
 Q_j=1+q+\cdots+q^j.
\]
Set $Q_0=1$. A point $a\in\PG(V_j)$ and a dual point
$b\in\PG(V_j^*)$ are incident, written $a\perp b$, if every covector
in $b$ vanishes on $a$. Thus a dual point also represents a hyperplane
in $\PG(V_j)$. All dimensions of linear spans below are \emph{vector}
dimensions; a vector subspace of dimension $r\ge1$ has $Q_{r-1}$
projective points. For sets in the opposite projective spaces, write
$e(S,T)$ for their number of incident pairs.

\begin{lemma}[Weighted incidence estimates]
\label{base:incidence}
Let $j\ge2$, and put $p_j=Q_{j-1}/Q_j=(1-Q_j^{-1})/q$.
For any real weights $w_x$ on one of the two projective spaces,
\begin{equation}
 \sum_y\left(\sum_{x\perp y}w_x-p_j\sum_xw_x\right)^2
 =q^{j-1}\left(\sum_xw_x^2-\frac{(\sum_xw_x)^2}{Q_j}\right)
 \le q^{j-1}\sum_xw_x^2.
 \label{base:weighted-variance}
\end{equation}
In particular, for sets $S,T$ in the opposite spaces,
\begin{equation}
 \left|e(S,T)-p_j|S||T|\right|
 \le q^{(j-1)/2}\sqrt{|S||T|}.
 \label{base:mixing}
\end{equation}
If $S,T$ are nonempty and $e(S,T)/(|S||T|)=o(1/q)$, then
$|S||T|\le Cq^{j+1}$ for all sufficiently large $q$.
\end{lemma}

\begin{proof}
Let $M$ be the point-hyperplane incidence matrix. Every point lies in
$Q_{j-1}$ hyperplanes, and two distinct points lie in $Q_{j-2}$
hyperplanes. Hence, writing $J$ for the all-ones matrix,
\[
 MM^{\mathsf T}=q^{j-1}I+Q_{j-2}J.
\]
The same identity holds with the two projective spaces interchanged.
Subtract the mean of $w$ and apply this identity to obtain
Equation~\eqref{base:weighted-variance}. Applying the resulting operator
norm bound to the centered indicators of $S$ and $T$ proves
Equation~\eqref{base:mixing}. Finally, $p_j=(1+o(1))/q$, so the sparsity
assumption and Equation~\eqref{base:mixing} give
\[
 \frac{|S||T|}{2q}
 \le q^{(j-1)/2}\sqrt{|S||T|}
\]
for large $q$, proving the last assertion.
\end{proof}

\subsection{The flag graph and rectangles}

We use Brada\v{c}'s ordered flag construction
\cite[Section~2.5]{Bradac2026}. A \emph{flag} in dimension $d$ is an incident pair $(a,b)$ in
$\PG(V_d)\times\PG(V_d^*)$. There are $Q_dQ_{d-1}$ flags. Sample an
ordered stream of $N$ independent uniform flags $(a_i,b_i)$, and make
a graph on its positions $1,\ldots,N$ by declaring, for $i<j$,
\begin{equation}
 ij\in E(G)\quad\Longleftrightarrow\quad
 a_i\perp b_j\ \hbox{ and }\ a_j\not\perp b_i.
 \label{base:edge-rule}
\end{equation}
Repeated flags are allowed: graph vertices are positions, not flag
values. An ordered flag tuple is called \emph{consistent} if
\begin{equation}
 i<j,\quad a_i\perp b_j\quad\Longrightarrow\quad a_j\perp b_i.
 \label{base:consistency}
\end{equation}

\begin{lemma}[Flag graph]
\label{base:flag-graph}
Every graph defined by Equation~\eqref{base:edge-rule} is $K_{d+1}$-free.
Its independent sets, in increasing position order, are exactly its
consistent subtuples. Reversing a consistent tuple and interchanging
its first and second endpoints preserves consistency.
\end{lemma}

\begin{proof}
Consider a clique with flags $(a_1,b_1),\ldots,(a_m,b_m)$ in its
increasing position order. For $r<m$, the covector $b_r$ annihilates
$a_1,\ldots,a_r$ but not $a_{r+1}$. The first assertion follows from
the earlier clique edges and the flag condition; the second follows
from the edge between positions $r$ and $r+1$. Thus $a_{r+1}$ lies
outside the vector span of its predecessors. The $m$ lines are
independent, while the nonzero covector $b_m$ annihilates all of them.
Its kernel has dimension $d$, so $m\le d$.
Negating Equation~\eqref{base:edge-rule} gives
Equation~\eqref{base:consistency}. Reversal and dual interchange turn
that same implication into itself.
\end{proof}

For a vector subspace $U\subseteq V_d^*$, define the rectangle
\[
 \mathcal R(U)=
 \{(a,b):b\subseteq U,\ a\subseteq U^\perp\}.
\]
All pairs in this set are flags. The next estimate holds simultaneously
for every rectangle, including a rectangle chosen from previously
observed data.

\begin{lemma}[Simultaneous rectangle occupancy]
\label{base:rectangles}
There is a constant $C_d$ such that, with probability $1-o(1)$, every
$\mathcal R(U)$ contains at most $C_d\sigma$ positions of the uniform
stream. The same event controls rectangles after dual interchange.
\end{lemma}

\begin{proof}
If $\dim U=r$ with $1\le r\le d$, then
\[
 |\mathcal R(U)|=Q_{r-1}Q_{d-r}\le4q^{d-1}.
\]
The rectangles for $r=0,d+1$ are empty. Thus the occupancy of any
fixed rectangle is binomial with mean at most $4\sigma$. There are
at most $(d+2)q^{(d+1)^2}$ subspaces, by counting possible ordered
bases. For $A>4e$, the exponential-moment bound for a binomial variable
gives
\[
 \Prob\{\text{occupancy}>A\sigma\}
 \le (4e/A)^{A\sigma}
 =q^{-A\log(A/(4e))},
\]
with an immaterial adjustment for rounding the threshold. Choose $A$
large enough in terms of $d$ and sum over the subspaces. Under dual
interchange, $\mathcal R(U)$ becomes the rectangle defined by
$U^\perp$, so no further event is needed.
\end{proof}

\subsection{Entropy after arbitrary selection}

For a discrete random variable $X$ with mass function $p$, write
$H(X)=-\sum_xp(x)\log p(x)$, with $0\log0=0$. Entropies use natural
units. Conditional entropy is the average of the corresponding
conditional entropies. Expanding these sums gives the chain rule and
\[
 I(X;Y)=H(X)+H(Y)-H(X,Y)
       =\KL(P_{X,Y}\Vert P_X\otimes P_Y),
\]
where $\KL(P\Vert Q)=\sum_xP(x)\log(P(x)/Q(x))$ is infinite if
$P$ is not absolutely continuous with respect to $Q$. These identities
also apply to countable variables of finite entropy.

\begin{lemma}[Selected-tuple entropy]
\label{base:selection-entropy}
Let a random stream have law dominated by $C_0$ times the iid uniform
flag-stream law, where $C_0$ is independent of $q$. From this stream,
using arbitrary dependent auxiliary randomness, extract a tuple $F$
of deterministic length $\ell=\Theta(k)$. Its order may be either
increasing stream order or reversed stream order with the endpoints
interchanged. Then
\begin{equation}
 H(F)\ge d\sigma\ell+\eta\ell\log\sigma-O(k).
 \label{base:entropy-lower}
\end{equation}
The implicit constant may depend on $C_0$ and the fixed comparison
constants in $\ell=\Theta(k)$.
\end{lemma}

\begin{proof}
For any prescribed tuple $f$, the event $F=f$ implies that $f$ occurs
on some increasing set of $\ell$ stream positions, in one of the two
allowed orientations. For each such set and orientation its iid
probability is $(Q_dQ_{d-1})^{-\ell}$. The stream domination and a
union bound therefore give
\[
 \Prob(F=f)\le2C_0\binom N\ell(Q_dQ_{d-1})^{-\ell}.
\]
This containment of events uses no assumption about the selector's
conditional distribution. Shannon entropy is at least minus the
logarithm of the largest atom. Now
\[
 \log(Q_dQ_{d-1})=(2d-1)\sigma+O(1/q),
 \qquad
 \log(eN/\ell)=(d-1)\sigma-\eta\log\sigma+O(1).
\]
Use $\binom N\ell\le(eN/\ell)^\ell$ and $\ell=\Theta(k)$ to conclude.
\end{proof}

The positive term $\eta\ell\log\sigma$ is the eventual contradiction:
we shall successively describe a selected tuple using only
$d\sigma\ell+O(k)$ entropy. Constant-probability restrictions of the
stream law are compatible with Lemma~\ref{base:selection-entropy}.
Indeed, if the current domination factor is $C_0$ and an event in the
joint probability space has probability at least $c>0$, then after
conditioning the stream marginal is dominated by $C_0/c$ times the
original law. This remains true when that event also uses auxiliary
randomness. A fixed number of such restrictions preserves a constant
domination factor.

\begin{lemma}[Transferring a rare event]
\label{base:event-transfer}
For two probability laws $P,Q$ on the same finite or countable space
and an event $E$,
\begin{equation}
 Q(E)\le\frac{\KL(P\Vert Q)+P(E)}{1-e^{-1}}.
 \label{base:kl-event}
\end{equation}
\end{lemma}

\begin{proof}
The assertion is immediate when the relative entropy is infinite.
Otherwise Jensen's inequality gives, for bounded $f$,
\[
 \KL(P\Vert Q)\ge \E_P f-\log\E_Q e^f.
\]
Apply this with $f=-1_E$. Since $-\log(1-u)\ge u$ for $0\le u<1$,
\[
 \KL(P\Vert Q)
 \ge-P(E)-\log\bigl(1-(1-e^{-1})Q(E)\bigr)
 \ge-P(E)+(1-e^{-1})Q(E),
\]
which proves the claim.
\end{proof}

\begin{lemma}[Elementary sampling tails]
\label{base:tails}
If $X$ is binomial or Poisson with mean $\mu$, then
$\E e^{tX}\le\exp((e^t-1)\mu)$ for every real $t$. For each fixed
$0<\delta<1$ there is $c_\delta>0$ such that
\[
 \Prob(|X-\mu|\ge\delta\mu)\le2e^{-c_\delta\mu}.
\]
For $h\ge e\mu>0$,
\[
 \Prob(X\ge h)\le(e\mu/h)^h.
\]
When $\mu=0$, the variable is zero almost surely.
\end{lemma}

\begin{proof}
For a Bernoulli variable of mean $p$, its exponential moment is
$1+p(e^t-1)\le e^{p(e^t-1)}$; multiply over independent trials.
For a Poisson variable the moment formula is an equality, by summing
its series. Exponential Markov with $t=\log(1+\delta)$ or
$t=\log(1-\delta)$ gives the two proportional tails. For the final
bound use $t=\log(h/\mu)$ and discard the factor $e^{-\mu}$.
\end{proof}

\subsection{The scales of a description}

The geometric descriptions and the compression stages use the following
common scales. Their input is a deterministic number $D$ measuring the
current description budget; the precise budget interpretation will be
given when a compression stage is defined. At present we need only its
allowed interval. Set
\begin{equation}
\begin{gathered}
 \beta=\eta/10^7,\qquad
 \sigma^\beta\le D\le\sigma^{1-\eta/2},\\
 K=D\sigma^{3\beta},\qquad K_*=D\sigma^{6\beta},\qquad
 L=D\sigma^{8\beta},\qquad P=LR,\\
 \sigma^\beta\le R\le2\sigma^\beta+2,
 \qquad R\text{ an even integer chosen deterministically}.
\end{gathered}
\label{base:parameters}
\end{equation}
The choice $R=2\lceil\sigma^\beta/2\rceil$ is one possibility.

\begin{lemma}[Separation of scales]
\label{base:scale-comparisons}
Uniformly over the interval for $D$ in Equation~\eqref{base:parameters},
\[
 K_*=o(L),\qquad L=o(P),\qquad P=o(\sigma),\qquad
 R\log L+\log\sigma=o(P).
\]
Moreover, for sufficiently large $q$,
$D\sigma^{9\beta}\le P\le4D\sigma^{9\beta}$.
\end{lemma}

\begin{proof}
The last bounds follow from the interval for $R$.
Also $K_*/L=\sigma^{-2\beta}$ and $L/P=1/R\le\sigma^{-\beta}$.
Since $9\beta<\eta/2$,
\[
 P\le4\sigma^{1-\eta/2+9\beta}=o(\sigma).
\]
Finally $L\ge\sigma^{9\beta}$, $P\ge\sigma^{10\beta}$, and
$\log L=O(\log\sigma)$, so
\[
 \frac{R\log L}{P}=\frac{\log L}{L}=o(1),
 \qquad \frac{\log\sigma}{P}=o(1).
\]
\end{proof}
